\section*{Chapter \ref{ch:b3:oligopoly-and-collusion} --- Oligopoly and Collusion}

\begin{solution}{exo:b3:oligopoly-and-collusion:1}
Against private label the six changes are 4, 8, 6, 8, 8 and 4~\%: median 7~\%. Against other brands they
are 0, 4, 2, 6, 4 and $-0.9$~\%: median 3~\%. One estimate is negative. % figdata: mergers
\end{solution}

\begin{solution}{exo:b3:oligopoly-and-collusion:2}
Prices $c + t = 14$; each firm serves half the consumers and earns $4 \times 0.5 = 2$.
\end{solution}

\begin{solution}{exo:b3:oligopoly-and-collusion:3}
$1/(0.5 \times 0.8) = 2.5$ dollars.
\end{solution}

\begin{solution}{exo:b3:oligopoly-and-collusion:4}
Prices that move closely with each firm's own costs, differ across firms in step
with their cost differences, and vary as much as in periods known to be competitive.
\end{solution}

\begin{solution}{exo:b3:oligopoly-and-collusion:5}
$1 - 1/2 = 0.5$ and $1 - 1/5 = 0.8$. With $\beta = 1/(1 + r)$, collusion among 5 firms holds while $1/(1 + r) \ge 0.8$, that is
$r \le 25~\%$ per period.
\end{solution}

\begin{solution}{exo:b3:oligopoly-and-collusion:6}
Before, each markup is $1/(1 - 0.25) = 1.33$. With the merged share 0.5 held fixed, the merged firm wants
$1/(1 - 0.5) = 2$. In equilibrium its higher prices lower its share, which lowers the markup
it wants, while the rival's higher price raises it a little; the equilibrium markup is
below 2.
\end{solution}

\begin{solution}{exo:b3:oligopoly-and-collusion:7}
$0.3 \times 4 = 1.2$ dollars. A saving $E$ with $E \times 10 = 1.2$, that is 12~\% of marginal cost.
\end{solution}

\begin{solution}{exo:b3:oligopoly-and-collusion:8}
The tested solution is \code{ex08_thresholds.py}, in the chapter's \code{exercises/}
folder. The thresholds per dentist are 0.53, 0.71, 0.83, 0.85 and 0.87 thousand residents: they
rise steeply with the second and third dentist and barely after, so entry beyond the
third changes conduct little. % output: ex08_thresholds
\end{solution}

\begin{solution}{exo:b3:oligopoly-and-collusion:9}
The rise would shrink or vanish when the variable input includes marketing and
management costs; it would be concentrated in firms whose share of those costs grew;
and demand-side markups or profit rates for the same firms would not rise with it.
\end{solution}

\begin{solution}{exo:b3:oligopoly-and-collusion:10}
The proof is in the text: the derivatives of the logit shares and a common markup
across the owner's products.
\end{solution}

\begin{solution}{exo:b3:oligopoly-and-collusion:11}
The tested solution is \code{ex11_pairs.py}, in the chapter's \code{exercises/}
folder. Merging with the close substitute, product 2, gives \omterm{def:b3:oligopoly-and-collusion:mergers}{upward pricing pressure} 0.568 and
a price rise of 0.460; merging with product 4 gives 0.100 and 0.079. The test ranks the mergers
correctly and, holding rivals' prices fixed and passing on less than all of the
pressure, overstates each rise. % output: ex11_pairs
\end{solution}

\begin{solution}{exo:b3:oligopoly-and-collusion:12}
If the other stays silent, staying silent costs $qF$ and reporting 0: reporting is better
for any $q > 0$. If the other reports, staying silent costs $F$ and reporting $F/2$: reporting is better.
Each member therefore expects the other to report, so the trust a \omterm{def:b3:oligopoly-and-collusion:cartel}{cartel} needs
cannot hold; the first to report gains most, which is why members race to the
authority.
\end{solution}

\begin{solution}{pb:b3:oligopoly-and-collusion:bakeries}
% figdata: bakeries_stats
\textbf{1.} $s_j = e^{\delta_j - 0.8p_j}/(1 + \sum_k e^{\delta_k - 0.8p_k})$.

\textbf{2.} $p_j - 2 = 1/(0.8(1 - s_j))$.

\textbf{3.} $p_1 = p_2 = 3.635$, $p_3 = 3.514$ dalens.

\textbf{4.} 0.2357, 0.2357 and 0.1741; no loaf: 0.3545.

\textbf{5.} $1/(0.8 \times (1 - 0.2357)) = 1.635 = 3.635 - 2$.

\textbf{6.} $1.635 \times 0.2357 \times 3{,}000$, about 1{,}156 dalens a day.

\textbf{7.} 3{,}889 dalens a day.

\textbf{8.} A loaf lost by bakery 1 is now partly recaptured by bakery 2, under the same
owner: raising either price costs the merged firm less than it cost each bakery alone.

\textbf{9.} $p_1 = p_2 = 4.055$, $p_3 = 3.555$.

\textbf{10.} 0.1959, 0.1959 and 0.1960; no loaf: 0.412.

\textbf{11.} The merged bakeries by 0.420 dalens; bakery 3 by 0.041.

\textbf{12.} $3{,}889 - 3{,}323 = 566$ dalens a day.

\textbf{13.} 2{,}415 dalens a day, against $2 \times 1{,}156 = 2{,}312$: a gain of 103.

\textbf{14.} Prices are strategic complements: its rivals' higher prices send it
customers and raise the markup it wants.

\textbf{15.} $D_{12} = 0.2357/(1 - 0.2357) = 0.308$.

\textbf{16.} 1.635 dalens.

\textbf{17.} $0.308 \times 1.635 = 0.504$ dalens a loaf.

\textbf{18.} $0.504/2 = 25.2~\%$ of marginal cost.

\textbf{19.} $1 - 1.27/2 = 36.5~\%$.

\textbf{20.} The test holds bakery 3's price fixed and asks only whether the price
wants to rise; in the simulation bakery 3 raises its price too, and a cost saving is
passed on to the price only in part, so a larger saving is needed.

\textbf{21.} $0.40 - (4.05 - 3.90) = 0.25$ dalens.

\textbf{22.} From 3.603 to 3.888: 0.286, about 0.29 dalens.

\textbf{23.} The measured rise is about seven-eighths of the predicted one ($0.25/0.286$).

\textbf{24.} Costs, such as flour or wages, may move differently in the two towns; and
the mix of loaves sold may change, so that the mean price moves without any price
changing.

\textbf{25.} Cost savings from the merger, entry or the threat of it, or conduct other
than the Bertrand pricing assumed.

\textbf{26.} The simulation predicts that the merged bakeries raise their prices by 0.42
dalens and the mean price rises by 0.29; the retrospective measures 0.25. A cost saving of 36.5~\% would
have made the merger neutral for the merged bakeries' prices (25.2~\% by the upward-pricing-pressure test).
\end{solution}
